\documentclass[a4paper,12pt]{article}
\usepackage[russian]{babel}
\usepackage[utf8]{inputenc}
\usepackage{amsmath, amssymb}
\usepackage{listings}
\usepackage{xcolor}
\usepackage{geometry}
\geometry{top=2cm, bottom=2cm, left=2.5cm, right=2.5cm}

\lstset{
    language=Haskell,
    basicstyle=\ttfamily\small,
    keywordstyle=\color{blue},
    commentstyle=\color{green!50!black},
    stringstyle=\color{red},
    numbers=left,
    numberstyle=\tiny\color{gray},
    stepnumber=1,
    breaklines=true,
    tabsize=2,
    showspaces=false,
    showstringspaces=false
}

\title{Интерпретатор Mini-Lisp на Haskell \\[0.3cm]
\large }
\author{}
\date{}

\begin{document}

\maketitle

\section{Формальная постановка задачи}

\subsection{Синтаксис языка}

\begin{verbatim}
<expr>      ::= <number> | <atom> | <list> | <quoted>

<number>    ::= '-'? <digit>+
<digit>     ::= '0'..'9'

<atom>      ::= <letter> <atom-rest>*
<atom-rest> ::= <letter> | <digit> | '-' | '_' | '?' | '!'
<letter>    ::= 'a'..'z' | 'A'..'Z' | '+' | '-' | '*' | '/' | '=' | '<' | '>'

<list>      ::= '(' <expr>* ')'                    (* обычный список *)
              | '(' <expr> '.' <expr> ')'          (* точечная пара *)

<quoted>    ::= '\'' <expr>

<expr>      ::= <number> | <atom> | <list> | <quoted>
\end{verbatim}

\subsection{Семантика}

\subsubsection{Типы данных}
\begin{itemize}
    \item \texttt{Atom String} — атом (символьное имя)
    \item \texttt{Number Integer} — целое число
    \item \texttt{List [Val]} — обычный список
    \item \texttt{DottedPair Val Val} — точечная пара
\end{itemize}

\subsubsection{Встроенные константы}
\begin{itemize}
    \item \texttt{t} — логическая истина
    \item \texttt{nil} или \texttt{()} — пустой список (логическая ложь)
\end{itemize}

\subsubsection{Специальные формы}
\begin{itemize}
    \item \texttt{(quote <expr>)} — возвращает выражение без вычисления
    \item \texttt{(if <cond> <then> <else>)} — условное выражение
    \item \texttt{(define <var> <expr>)} — определение переменной
\end{itemize}

\subsubsection{Встроенные функции}
\begin{itemize}
    \item \texttt{(car <list>)} — первый элемент списка или точечной пары
    \item \texttt{(cdr <list>)} — хвост списка или второй элемент пары
    \item \texttt{(cons <x> <y>)} — конструирование пары или списка
    \item \texttt{(eq <x1> <x2>)} — сравнение на равенство
    \item \texttt{(+ <x1> <x2>)}, \texttt{(- <x1> <x2>)}, \texttt{(* <x1> <x2>)}, \texttt{(/ <x1> <x2>)} — арифметические операции
    \item \texttt{(< <x1> <x2>)}, \texttt{(> <x1> <x2>)}, \texttt{(= <x1> <x2>)} — сравнение чисел
\end{itemize}

\section{Описание алгоритмов}

\subsection{Лексический анализ и парсинг}

Для разбора используется библиотека \texttt{Parsec} — реализация комбинаторных парсеров.

\subsubsection{Алгоритм парсинга чисел}
\begin{enumerate}
    \item Опционально обрабатываем знак минуса (только если после него идёт цифра)
    \item Парсим одну или более цифр
    \item Преобразуем строку в целое число
    \item Применяем знак (если был)
\end{enumerate}

\subsubsection{Алгоритм парсинга атомов}
\begin{enumerate}
    \item Первый символ: буква или один из символов \texttt{+-*/=<>}
    \item Последующие символы: буквы, цифры или \texttt{-\_?!}
    \item Формируем атом с полученным именем
\end{enumerate}

\subsubsection{Алгоритм парсинга списков}
Парсер использует стратегию с возвратом (\texttt{try}):

\begin{enumerate}
    \item Пробуем распознать пустой список \texttt{()}
    \item Если не пустой — пробуем распознать точечную пару \texttt{(a . b)}
    \item Иначе — распознаём обычный список \texttt{(a b c ...)}
\end{enumerate}

Для обычного списка:
\begin{enumerate}
    \item Читаем открывающую скобку
    \item Парсим первое выражение
    \item Повторяем парсинг следующих выражений (ноль или более)
    \item Читаем закрывающую скобку
\end{enumerate}

\subsubsection{Алгоритм парсинга quoted-выражений}
\begin{enumerate}
    \item Распознаём символ \texttt{'}
    \item Парсим следующее выражение
    \item Преобразуем в \texttt{(quote <expr>)}
\end{enumerate}

\subsection{Интерпретация}

\subsubsection{Управление переменными}
\begin{itemize}
    \item \textbf{getVar} — линейный поиск переменной в списке окружения (ассоциативном списке)
    \item \textbf{setVar} — добавление новой пары (имя, значение) в начало окружения (реализует динамическую область видимости)
    \item Запрещено переопределение встроенных констант \texttt{t}, \texttt{nil}, \texttt{T}, \texttt{NIL}
\end{itemize}

\subsubsection{Приведение к логическому значению}
Функция \texttt{toBool}:
\begin{itemize}
    \item Пустой список \texttt{(List [])} → \texttt{False}
    \item Атом \texttt{"nil"} → \texttt{False}
    \item Атом \texttt{"t"} → \texttt{True}
    \item Любое другое значение → \texttt{False}
\end{itemize}

\subsubsection{Алгоритм eval (вычисление выражения)}
Функция \texttt{eval :: Variables -> Val -> Either String (Val, Variables)}:

\begin{enumerate}
    \item \textbf{Число} — возвращается без изменений
    \item \textbf{Атом (переменная)} — ищется в окружении; если не найдена — ошибка
    \item \textbf{Список} — анализируется первый элемент:
    \begin{itemize}
        \item \texttt{quote} — возвращает второй аргумент без вычисления
        \item \texttt{if} — вычисляется условие; в зависимости от истинности вычисляется одна из ветвей
        \item \texttt{define} — вычисляется выражение, переменная добавляется в окружение
        \item \texttt{car}, \texttt{cdr}, \texttt{cons}, \texttt{eq} — вычисляются аргументы, затем выполняется операция
        \item \texttt{+}, \texttt{-}, \texttt{*}, \texttt{/}, \texttt{<}, \texttt{>}, \texttt{=} — вычисляются оба аргумента (с последовательным обновлением окружения), преобразуются к числам, выполняется операция
        \item \textit{иначе} — ошибка (в данной реализации функции не могут быть определены пользователем)
    \end{itemize}
    \item \textbf{Пустой список} — возвращается пустой список
\end{enumerate}

\subsubsection{Алгоритм работы с точечными парами}
\begin{itemize}
    \item \texttt{car} от \texttt{DottedPair x y} → \texttt{x}
    \item \texttt{cdr} от \texttt{DottedPair x y} → \texttt{y}
    \item \texttt{cons}:
    \begin{itemize}
        \item Если второй аргумент — \texttt{List ys}, результат — \texttt{List (x:ys)}
        \item Иначе — \texttt{DottedPair x y}
    \end{itemize}
\end{itemize}

\subsection{REPL}

Алгоритм работы REPL:
\begin{enumerate}
    \item Вывод \texttt{lisp>}
    \item Чтение строки ввода
    \item Проверка на команду \texttt{exit}
    \item Парсинг строки в \texttt{Val}
    \item Вычисление выражения в текущем окружении
    \item Вывод результата
    \item Обновление окружения (при \texttt{define})
    \item Повторение цикла
\end{enumerate}

\subsection{Обработка ошибок}
Все ошибки обрабатываются с использованием типа \texttt{Either String a}:
\begin{itemize}
    \item Ошибки парсинга — сообщения от Parsec
    \item Ошибки выполнения — сообщения с описанием проблемы:
    \begin{itemize}
        \item Неопределённая переменная
        \item Ожидание числа
        \item Деление на ноль
        \item Переопределение встроенной константы
        \item Некорректный аргумент для \texttt{car}/\texttt{cdr}
    \end{itemize}
\end{itemize}

\section{Структура исходного кода}

Исходный код организован следующим образом:

\begin{lstlisting}
import Text.Parsec
import Text.Parsec.String (Parser)

data Val = Atom String | Number Integer | List [Val] | DottedPair Val Val
  deriving (Show, Eq)

type Variables = [(String, Val)]

-- === Парсер ===
whiteSpace :: Parser ()
atomParser :: Parser Val
numberParser :: Parser Val
emptyListParser :: Parser Val
dottedPairParser :: Parser Val
regularListParser :: Parser Val
listParser :: Parser Val
quotedParser :: Parser Val
parseExpr :: Parser Val
parseLisp :: String -> Either ParseError Val

-- === Управление переменными ===
isForbiddenName :: String -> Bool
getVar :: String -> Variables -> Maybe Val
setVar :: String -> Val -> Variables -> Either String Variables

-- === Вспомогательные функции ===
toNumber :: Val -> Either String Integer
toBool :: Val -> Bool

-- === Интерпретатор ===
eval :: Variables -> Val -> Either String (Val, Variables)

-- === REPL ===
repl :: Variables -> IO ()
initialEnv :: Variables
main :: IO ()
\end{lstlisting}

\subsection{Примечания по реализации}
\begin{itemize}
    \item Парсеры используют комбинатор \texttt{try} для реализации возврата при разборе альтернатив
    \item Пробелы игнорируются с помощью \texttt{whiteSpace}
    \item Окружение передаётся явно через аргумент функции \texttt{eval}
    \item Каждая операция возвращает новое окружение на случай изменений
    \item При бинарных операциях окружение последовательно обновляется при вычислении каждого аргумента
\end{itemize}

\end{document}